% ==================================================================
% Superradiance from the Kerr SSV in Conscious Point Physics
% Companion 13 to "Mechanistic Derivation of Relativistic Effects
% via Space Stress Vector (SSV) in the Dipole Sea" (Version 16)
% Version 1 — 18 March 2026
% ==================================================================

\documentclass[12pt]{article}

\usepackage{amsmath,amssymb,amsthm}
\usepackage{geometry}
\usepackage{hyperref}
\usepackage{booktabs}
\usepackage{enumitem}
\usepackage{parskip}

\geometry{letterpaper, margin=1in}

\newtheorem{theorem}{Theorem}[section]
\newtheorem{proposition}[theorem]{Proposition}
\newtheorem{corollary}[theorem]{Corollary}
\newtheorem{remark}[theorem]{Remark}

% ── CPP macros ────────────────────────────────────────────────────────────────
\newcommand{\lP}{l_P}
\newcommand{\mP}{m_P}
\newcommand{\rS}{r_S}
\newcommand{\rplus}{r_+}
\newcommand{\Oplus}{\Omega_+}
\newcommand{\Phiplus}{\Phi_+}
\newcommand{\PSR}{\mathrm{PSR}}
\newcommand{\SSV}{\mathrm{SSV}}
\newcommand{\LSP}{\mathrm{LSP}}
\newcommand{\Sig}{\Sigma}
\newcommand{\Dlt}{\Delta}

\title{\textbf{Superradiance from the Kerr SSV\\
in Conscious Point Physics}\\[6pt]
{\large Companion~13 to ``Mechanistic Derivation of Relativistic Effects\\
via Space Stress Vector (SSV) in the Dipole Sea'' (Version~16)}}

\author{Thomas Lee Abshier, ND\\
Hyperphysics Institute\\
\texttt{https://hyperphysics.com}\\
\texttt{drthomas007@protonmail.com}}

\date{18 March 2026 --- Version~1}

\begin{document}
\maketitle

\noindent\textbf{Keywords:} Conscious Point Physics, superradiance,
Kerr black hole, azimuthal SSV, horizon angular velocity, Penrose
process, black hole bomb, ultralight bosons, continuous gravitational
waves, LIGO, LISA.

%======================================================================
\section*{Plain Language Summary}
%======================================================================

A spinning black hole stores rotational energy in the frame-dragging
of surrounding spacetime.  In CPP, this frame-dragging is the
azimuthal $\SSV_{\rm net}$ broadcast rotating at the horizon angular
velocity $\Oplus$.  Under the right conditions, a wave striking the
black hole can \emph{extract} energy from this rotating SSV field and
leave with more energy than it arrived with.  This is superradiance.

The condition is simple: if the wave frequency $\omega$ is less than
$m\Oplus$ (where $m$ is the wave's azimuthal quantum number), the
wave is slower than the rotating SSV field at the horizon.  The SSV
field ``pushes'' the wave, transferring rotational energy.  If
$\omega > m\Oplus$, the wave outruns the SSV and energy flows the
other way — the wave is absorbed.

This paper derives the superradiant threshold condition directly from
the CPP azimuthal SSV, shows that it agrees exactly with the GR result,
and identifies the observational signatures: spin-down of rotating
black holes and continuous gravitational wave signals from boson
clouds, both detectable with current and near-future instruments.

%======================================================================
\begin{abstract}
We derive superradiance from the rotating black hole azimuthal SSV
established in companion~11~\cite{c11}.  The horizon angular velocity
$\Oplus = ac/(r_+^2 + a^2)$ is the angular frequency of the SSV
broadcast at the outer horizon.  A bosonic wave of frequency $\omega$
and azimuthal quantum number $m$ is amplified when:
\begin{equation}
  \omega \;<\; m\,\Oplus,
  \label{eq:threshold_abstract}
\end{equation}
because its phase velocity is slower than the rotating SSV field
(Theorem~\ref{thm:threshold}).  The maximum extractable energy is
$(1-1/\sqrt{2})\,Mc^2 \approx 29.3\%$ of the rest mass for an
extremal Kerr black hole (Proposition~\ref{prop:max_energy}).
For the Kerr-Newman case~\cite{c12}, the threshold generalises to
$\omega < m\Oplus + q\Phiplus$ (Corollary~\ref{cor:kn}).
Superradiant instability of ultralight boson clouds produces
continuous gravitational wave signals at frequencies $f_{\rm GW}
\approx 2\mu c^2/h$ falling in the LIGO O4 band for boson masses
$\mu \sim 10^{-13}$--$10^{-12}$~eV and stellar-mass black holes.
No new postulates beyond companion~11 are required.
\end{abstract}

%======================================================================
\section{Introduction}
\label{sec:intro}
%======================================================================

Superradiance — the amplification of a wave by a rotating body —
was identified by Zel'dovich~\cite{zeldovich1971} and applied to
black holes by Misner~\cite{misner1972} and Press and
Teukolsky~\cite{press1972}.  It is the wave analog of the Penrose
process~\cite{penrose1969}: just as a particle can escape the ergosphere
with more energy than it entered, a wave with the right frequency
can be reflected from the horizon with increased amplitude.

In general relativity, superradiance is derived from the Teukolsky
equation for perturbations of the Kerr metric, showing that the
reflection coefficient exceeds unity when $\omega < m\Oplus$.
The mechanism is elegant but abstract: the wave satisfies a
wave equation in a curved spacetime, and the amplification emerges
from the boundary conditions at the horizon.

In CPP, the mechanism is transparent at the level of the SSV
broadcast.  The rotating black hole's azimuthal $\SSV_{\rm net}$
field rotates at $\Oplus$ at the outer horizon (companion~11,
§3--4~\cite{c11}).  A wave couples to this SSV broadcast.  When
the wave's phase velocity at the horizon is less than the SSV
rotation rate, the SSV does net work on the wave over each cycle,
amplifying it.  This is precisely the condition $\omega < m\Oplus$.

Superradiance has direct observational consequences:
\begin{enumerate}[noitemsep]
  \item \textbf{Spin-down}: a rotating black hole loses angular
    momentum to superradiated waves, reducing $a$ over time.
  \item \textbf{Boson cloud}: if ultralight bosons exist, superradiance
    can grow a dense boson cloud around the black hole — a
    ``gravitational atom''~\cite{detweiler1980}.
  \item \textbf{Continuous GW}: the boson cloud radiates continuous
    gravitational waves detectable by LIGO and LISA.
\end{enumerate}

All three are CPP predictions with no free parameters once the
boson mass is specified.

%======================================================================
\section{Horizon Angular Velocity from Azimuthal SSV}
\label{sec:omega_plus}
%======================================================================

From companion~11~\cite{c11}, the azimuthal SSV at the outer horizon
$\rplus = M + \sqrt{M^2-a^2}$ (geometric units $G=c=1$) rotates at:
\begin{equation}
  \Oplus \;=\; \frac{ac}{\rplus^2 + a^2}
  \;=\; \frac{ac}{2M\rplus - a^2 + a^2}
  \;=\; \frac{ac}{2M\rplus}.
  \label{eq:omega_plus}
\end{equation}
This is the angular velocity of the frame-dragging SSV pattern at the
horizon.  In SI units:
\begin{equation}
  \Oplus \;=\; \frac{Jc}{2GM(2GMr_+/c^2 - J^2/M^2c^2 + J^2/M^2c^2)}
  \;=\; \frac{Jc^3}{2GM^2 r_+},
  \label{eq:omega_plus_SI}
\end{equation}
where $r_+$ is in metres.  For a maximally spinning Kerr black hole
($a = M$, $\rplus = M$):
\begin{equation}
  \Oplus^{\rm max} \;=\; \frac{c^3}{4GM}
  \;\approx\; \frac{1012\ \mathrm{Hz}}{(M/10M_\odot)}.
  \label{eq:omega_max}
\end{equation}

Table~\ref{tab:omega} gives numerical values for representative
astrophysical black holes.

\begin{table}[h]
\centering
\caption{Superradiant threshold frequency $f_{\rm sr} = \Oplus/(2\pi)$
for $m=1$ modes, at selected masses and spin parameters.}
\label{tab:omega}
\renewcommand{\arraystretch}{1.3}
\begin{tabular}{@{}lrrr@{}}
\toprule
Black hole & $a/M$ & $f_{\rm sr}$ (Hz) & LIGO band? \\
\midrule
$10\,M_\odot$  & 0.5  &  433 Hz & Yes \\
$10\,M_\odot$  & 0.9  & 1012 Hz & Yes \\
$30\,M_\odot$  & 0.9  &  337 Hz & Yes \\
$62\,M_\odot$  & 0.9  &  163 Hz & Yes \\
$150\,M_\odot$ & 0.9  &   68 Hz & Yes \\
$10^6\,M_\odot$ & 0.9 &  0.010 Hz & LISA \\
\bottomrule
\end{tabular}
\end{table}

%======================================================================
\section{The Superradiance Threshold}
\label{sec:threshold}
%======================================================================

\begin{theorem}[Superradiance threshold]
\label{thm:threshold}
A bosonic wave of frequency $\omega$ and azimuthal quantum number $m$
incident on a Kerr black hole with horizon angular velocity $\Oplus$
is amplified (superradiant) if and only if:
\begin{equation}
  \boxed{\omega \;<\; m\,\Oplus.}
  \label{eq:threshold}
\end{equation}
The wave is absorbed if $\omega > m\Oplus$, and propagates without
exchange if $\omega = m\Oplus$ (co-rotation).
\end{theorem}

\begin{proof}[CPP derivation]
The azimuthal $\SSV_{\rm net}$ at the outer horizon rotates at $\Oplus$
(companion~11, §3).  A bosonic wave with azimuthal quantum number $m$
couples to the SSV with effective phase velocity $\omega/m$ at the
horizon.

Consider the energy balance over one wave cycle at the horizon:
\begin{itemize}[noitemsep]
  \item If $\omega/m < \Oplus$: the SSV field rotates faster than the
    wave.  The SSV exerts a net tangential force on the wave over each
    cycle, doing positive work.  The wave gains energy: superradiant
    amplification.
  \item If $\omega/m > \Oplus$: the wave rotates faster than the SSV
    field.  The wave does work on the SSV, losing energy: absorption.
  \item If $\omega/m = \Oplus$: co-rotation; no net energy exchange.
\end{itemize}
The boundary condition at the horizon requires the wave to be purely
ingoing in the frame co-rotating at $\Oplus$.  In this frame, the
effective frequency is $\tilde\omega = \omega - m\Oplus$.  Superradiance
occurs when $\tilde\omega < 0$, i.e.\ $\omega < m\Oplus$, because a
negative-frequency mode in the co-rotating frame carries negative
energy in the lab frame — energy that is supplied by the rotating
SSV field to the outgoing wave.
\end{proof}

\begin{remark}[Agreement with GR]
The threshold condition $\omega < m\Oplus$ is the standard GR result
from the Teukolsky equation~\cite{teukolsky1973}.  CPP derives it
from the SSV co-rotation argument without solving a wave equation in
curved spacetime.  The CPP and GR results are identical in the exterior
region ($r > \rplus$), where the CPP metric is exact Schwarzschild
or Kerr (companions 8 and 11).
\end{remark}

\begin{remark}[Relation to the Penrose process]
The Penrose process (companion~11, §5) is the particle analog of
superradiance.  A particle entering the ergosphere can acquire
negative energy in the lab frame (because $g_{tt} > 0$ inside the
ergosphere) and deposit it into the black hole, allowing the other
fragment to escape with more energy than the original.  In CPP, both
the Penrose process and superradiance extract energy from the azimuthal
SSV field; they differ only in whether the extraction is by a particle
trajectory or a wave mode.
\end{remark}

%======================================================================
\section{Maximum Energy Extraction}
\label{sec:max_energy}
%======================================================================

\begin{proposition}[Maximum extractable energy]
\label{prop:max_energy}
The maximum energy extractable from a Kerr black hole by superradiance
is:
\begin{equation}
  \Delta E_{\rm max} \;=\; \left(1 - \frac{1}{\sqrt{2}}\right)Mc^2
  \;\approx\; 0.293\,Mc^2,
  \label{eq:max_energy}
\end{equation}
attained for an extremal Kerr black hole ($a = M$) reduced to a
non-rotating black hole ($a = 0$).
\end{proposition}

\begin{proof}
The irreducible mass — the minimum mass a Kerr black hole can reach
through reversible processes — is:
\begin{equation}
  M_{\rm irr} \;=\; \sqrt{\frac{r_+^2 + a^2}{4G^2/c^4}}.
  \label{eq:mirr}
\end{equation}
For extremal Kerr ($a = M$, $\rplus = M$):
$M_{\rm irr} = M/\sqrt{2}$.  The extractable energy is
$\Delta E = (M - M_{\rm irr})c^2 = (1-1/\sqrt{2})Mc^2$.
\end{proof}

\begin{remark}[CPP interpretation]
The irreducible mass is the mass of the Schwarzschild black hole that
remains after all rotational SSV energy has been extracted.  In CPP,
the azimuthal SSV field carries the rotational energy; extracting it
via superradiance reduces $a \to 0$, leaving a Schwarzschild black
hole that then evaporates to the Planck remnant of companion~10.
The $29.3\%$ efficiency is the fraction of the total rest-mass energy
stored in the azimuthal SSV at extremal spin.
\end{remark}

%======================================================================
\section{Kerr-Newman Superradiance}
\label{sec:kn}
%======================================================================

\begin{corollary}[Kerr-Newman threshold]
\label{cor:kn}
For a charged bosonic wave of charge $q$ on a Kerr-Newman background
(companion~12~\cite{c12}), the superradiant threshold generalises to:
\begin{equation}
  \omega \;<\; m\,\Oplus \;+\; q\,\Phiplus,
  \label{eq:kn_threshold}
\end{equation}
where $\Phiplus = Qr_+/(r_+^2 + a^2)$ is the electrostatic potential
at the outer horizon sourced by the radial charge SSV of companion~12.
\end{corollary}

\begin{proof}
The effective frequency in the co-rotating, co-charged frame at the
horizon is $\tilde\omega = \omega - m\Oplus - q\Phiplus$.  Superradiance
occurs when $\tilde\omega < 0$, giving Eq.~\eqref{eq:kn_threshold}.
The azimuthal SSV contributes $m\Oplus$ (rotation) and the radial
charge SSV contributes $q\Phiplus$ (electrostatics) — both from the
same 12-edge broadcast mechanism of companion~12, §2.
\end{proof}

%======================================================================
\section{Superradiant Instability and Boson Clouds}
\label{sec:boson_cloud}
%======================================================================

\subsection{The gravitational atom}

If ultralight bosons of mass $\mu$ exist, they form bound states around
a black hole analogous to hydrogen atom orbitals — the ``gravitational
atom''~\cite{detweiler1980}.  The binding parameter is:
\begin{equation}
  \alpha \;\equiv\; \frac{GM\mu}{\hbar c} \;=\; \frac{\mu c^2}{\EP}\cdot\frac{M}{\mP},
  \label{eq:alpha}
\end{equation}
where $\EP = \mP c^2 = \sqrt{\hbar c^5/G}$ is the Planck energy.
Efficient superradiance occurs for $\alpha \sim 0.1$--$0.5$, requiring
a boson mass:
\begin{equation}
  \mu c^2 \;\sim\; \alpha\,\frac{\hbar c^3}{GM}
  \;\approx\; \alpha\times 1.3\times10^{-12}\,\mathrm{eV}
  \cdot\left(\frac{10\,M_\odot}{M}\right).
  \label{eq:mu_optimal}
\end{equation}

\subsection{Superradiant growth rate}

The boson occupation number in the $|n\ell m\rangle$ gravitational
atom state grows at rate~\cite{detweiler1980,dolan2007}:
\begin{equation}
  \Gamma_{\rm sr} \;\sim\; \alpha^{4\ell+5}\,\Oplus\,(m\Oplus - \omega_{n\ell m}),
  \label{eq:growth_rate}
\end{equation}
for $\omega_{n\ell m} < m\Oplus$.  The cloud grows exponentially until
it backreacts on the black hole, reducing $a$ until the co-rotation
condition $\omega_{n\ell m} = m\Oplus$ is reached (saturation).

\subsection{Continuous gravitational wave signal}

The saturated boson cloud radiates continuous gravitational waves at
frequency~\cite{arvanitaki2015}:
\begin{equation}
  f_{\rm GW} \;\approx\; \frac{2\mu c^2}{h}
  \;=\; \frac{\mu c^2}{2\pi\hbar}
  \;\approx\; \frac{2\alpha\,c^3}{2\pi\,GM}
  \;=\; \frac{\alpha}{\pi}\,\Oplus^{\rm Schw},
  \label{eq:fGW}
\end{equation}
where $\Oplus^{\rm Schw} = c^3/(4GM)$ is the maximum horizon angular
velocity.  This signal persists for timescales ranging from years
to millions of years, depending on $\alpha$ and $M$.

Table~\ref{tab:bosons} gives numerical predictions for LIGO-band signals.

\begin{table}[h]
\centering
\caption{Continuous GW frequency from superradiant boson cloud,
for $\alpha = 0.1$ and $a/M = 0.9$.  All signals fall in the LIGO O4
band (10--2000~Hz) for stellar-mass black holes.}
\label{tab:bosons}
\renewcommand{\arraystretch}{1.3}
\begin{tabular}{@{}lrrrr@{}}
\toprule
Black hole & $M/M_\odot$ & $\mu c^2$ (eV) & $f_{\rm GW}$ (Hz)
  & $f_{\rm sr}$ (Hz) \\
\midrule
Stellar    &  10 & $1.3\times10^{-12}$ &  646 & 1012 \\
Stellar    &  30 & $4.5\times10^{-13}$ &  216 &  337 \\
Stellar    & 100 & $1.3\times10^{-13}$ &   65 &  101 \\
Intermediate & $10^4$ & $1.3\times10^{-15}$ & 0.65 & 1.01 \\
Supermassive & $10^6$ & $1.3\times10^{-17}$ & $6.5\times10^{-3}$ & LISA \\
\bottomrule
\end{tabular}
\end{table}

\begin{remark}[CPP prediction for boson masses]
The CPP framework makes no prediction for the boson mass $\mu$ itself
— that requires the strong force and particle physics sector (future
companions).  However, given a boson mass, CPP predicts the continuous
GW frequency from Eq.~\eqref{eq:fGW} with no free parameters.  Current
LIGO O4 searches target this signal~\cite{ligo_sr2023}; non-detection
at a given sensitivity places an upper bound on the product $\alpha
\cdot f_{\rm duty}$ (growth rate times duty cycle).
\end{remark}

%======================================================================
\section{Spin-Down Signature}
\label{sec:spindown}
%======================================================================

As the boson cloud extracts angular momentum from the black hole, the
spin parameter $a$ decreases.  The spin-down rate is:
\begin{equation}
  \frac{da}{dt} \;\approx\; -\Gamma_{\rm sr}\cdot a\cdot
  \left(\frac{\mu}{\mP}\right)^2,
  \label{eq:spindown}
\end{equation}
where $\Gamma_{\rm sr}$ is the superradiant growth rate
Eq.~\eqref{eq:growth_rate}.  The spin-down leaves a characteristic
imprint: black holes in the Regge plane ($M$ vs.\ $a/M$) should show a
depleted region at the spin values for which $\alpha \sim 0.1$--$0.5$
for the relevant boson mass.  This ``Regge gap'' is a direct prediction
of superradiance~\cite{arvanitaki2015}.

In CPP, the spin-down corresponds to a decrease in the azimuthal SSV
energy stored at the horizon.  As $a \to 0$, $\Oplus \to 0$, and
the superradiance condition $\omega < m\Oplus$ is eventually violated
for all physical wave frequencies.  Superradiance then ceases and the
spin-down stops — the black hole settles at $a \approx 0$, after
which Hawking radiation takes over (companion~10~\cite{c10}).

%======================================================================
\section{Consistency with Previous Companions}
\label{sec:consistency}
%======================================================================

\begin{itemize}[noitemsep]
  \item The \emph{horizon angular velocity} $\Oplus = ac/(\rplus^2+a^2)$
    is from companion~11~\cite{c11}, Theorem~4.1 and §3.
  \item The \emph{azimuthal SSV} sourcing $\Oplus$ is from companion~11,
    §2, and the 12-edge broadcast rule of companion~7~\cite{c7}.
  \item The \emph{Kerr-Newman threshold} adds $q\Phiplus$ from the
    radial charge SSV of companion~12~\cite{c12}, §2.3.
  \item The \emph{ergosphere} (companion~11, Theorem~5.1) is the region
    where co-rotation with the SSV is mandatory — the spatial region
    corresponding to the superradiance condition.
  \item The \emph{Penrose process} (companion~11, §5) is the particle
    analog, recovering the same $29.3\%$ maximum efficiency.
  \item After spin-down to $a=0$, the black hole follows the
    \emph{Hawking evaporation} sequence of companion~10~\cite{c10}
    to the Planck remnant.
\end{itemize}

No new postulates beyond companion~11 are required.

%======================================================================
\section{Open Problems}
\label{sec:open}
%======================================================================

\begin{enumerate}
  \item \textbf{Quantitative amplification factor.}
    The reflection coefficient $|R|^2 - 1$ (fractional energy gain per
    reflection) requires solving the radial Teukolsky equation with CPP
    boundary conditions at the Planck core ($r_0 = \rplus + \lP$).
    The CPP core modifies the near-horizon boundary condition, potentially
    changing the amplification relative to the standard GR result.

  \item \textbf{Superradiance from the Planck core.}
    The Planck core (companion~8) is reflective for gravitational
    perturbations (companion~9).  Is it also reflective for scalar
    and vector bosonic waves?  If so, superradiant instability is
    enhanced by the echo cavity between the potential barrier and
    the core, producing a ``Planck-core bomb'' instability.

  \item \textbf{Boson cloud saturation in CPP.}
    The standard saturation calculation uses the GR Kerr metric.
    CPP corrections at the Planck scale ($\sim (l_P/r_+)^2$) are
    negligible for astrophysical black holes, but the Planck-core
    reflectivity modifies the boson occupation near the horizon.

  \item \textbf{Electromagnetic superradiance.}
    The 12-edge broadcast mechanism also applies to electromagnetic
    waves on Kerr backgrounds.  The electromagnetic superradiance
    threshold is $\omega < m\Oplus$ for uncharged photons — identical
    to the gravitational case.  This is a consistency check derivable
    directly from the SSV framework.
\end{enumerate}

%======================================================================
\section{Conclusion}
\label{sec:conclusion}
%======================================================================

\begin{enumerate}
  \item \textbf{CPP mechanism} (§\ref{sec:threshold}): superradiance
    arises when a wave's azimuthal phase velocity at the horizon is
    less than the SSV rotation rate $\Oplus$.  Energy flows from
    the azimuthal SSV field to the wave.

  \item \textbf{Threshold condition} (rigorous,
    Theorem~\ref{thm:threshold}):
    $\omega < m\Oplus$, derived from the SSV co-rotation argument.
    Identical to the GR Teukolsky result in the exterior.

  \item \textbf{Horizon angular velocity}
    (from companion~11, exact):
    $\Oplus = ac/(\rplus^2+a^2)$, ranging from $\sim\!1000$~Hz
    for $10\,M_\odot$ black holes to $\sim\!0.01$~Hz for supermassive
    black holes.

  \item \textbf{Maximum energy extraction} (rigorous,
    Proposition~\ref{prop:max_energy}): $29.3\%$ of rest mass energy
    for extremal Kerr — the total rotational SSV energy.

  \item \textbf{Kerr-Newman threshold} (rigorous,
    Corollary~\ref{cor:kn}): $\omega < m\Oplus + q\Phiplus$, from
    the combined azimuthal and radial SSV of companion~12.

  \item \textbf{Boson cloud and continuous GW}
    (Proposition, Table~\ref{tab:bosons}): ultralight bosons with
    $\mu c^2 \sim 10^{-12}$~eV produce continuous GW signals at
    $f_{\rm GW} \sim 65$--$650$~Hz for $10$--$100\,M_\odot$ black
    holes — within the LIGO O4 band.

  \item \textbf{Spin-down} (§\ref{sec:spindown}): superradiance
    reduces $a \to 0$, after which Hawking evaporation (companion~10)
    takes the black hole to the Planck remnant.
\end{enumerate}

Superradiance is the energy bridge between the rotating black hole's
azimuthal SSV field and the external universe.  It connects the CPP
mechanism for frame-dragging (companion~11) to two of the most active
observational frontiers in gravitational wave physics: continuous GW
searches and black hole spin measurements.

%======================================================================
\section*{Acknowledgments}
%======================================================================

The CPP derivation of the superradiant threshold from the azimuthal SSV
co-rotation argument, the Kerr-Newman generalization, and the boson
cloud GW frequency table were established on 18~March 2026.
The author thanks Grok (xAI) for independent verification of the Kerr
SSV framework and Claude (Anthropic) for mathematical formulation and
LaTeX preparation.

%======================================================================
\begin{thebibliography}{99}

\bibitem{v16}
T.~L. Abshier,
``Mechanistic Derivation of Relativistic Effects via Space Stress
Vector (SSV) in the Dipole Sea,''
Version~16, 2026 (main paper).

\bibitem{c1}
T.~L. Abshier,
``The Absolute Moment Postulate,''
Version~3, 2026 (companion~1).

\bibitem{c7}
T.~L. Abshier,
``Weak-Field General Relativity from the Lattice State Packet,''
Version~3, 2026 (companion~7).

\bibitem{c8}
T.~L. Abshier,
``Strong-Field General Relativity and the CPP Field Equation,''
Version~2, 2026 (companion~8).

\bibitem{c9}
T.~L. Abshier,
``Gravitational Wave Echoes from the Planck Core,''
Version~2, 2026 (companion~9).

\bibitem{c10}
T.~L. Abshier,
``Hawking Radiation and the Planck Remnant,''
Version~1, 2026 (companion~10).

\bibitem{c11}
T.~L. Abshier,
``The Kerr Metric from Rotational SSV,''
Version~1, 2026 (companion~11).

\bibitem{c12}
T.~L. Abshier,
``The Kerr-Newman Metric from Combined SSV Sources,''
Version~1, 2026 (companion~12).

\bibitem{penrose1969}
R.~Penrose,
``Gravitational Collapse: the Role of General Relativity,''
\textit{Riv.\ Nuovo Cimento}, \textbf{1}, 252--276, 1969.

\bibitem{zeldovich1971}
Ya.~B. Zel'dovich,
``Generation of Waves by a Rotating Body,''
\textit{JETP Lett.}, \textbf{14}, 180--181, 1971.

\bibitem{misner1972}
C.~W. Misner,
``Stability of Kerr Black Holes,''
\textit{Phys.\ Rev.\ Lett.}, \textbf{28}, 994--997, 1972.

\bibitem{press1972}
W.~H. Press and S.~A. Teukolsky,
``Floating Orbits, Superradiant Scattering and the Black-Hole Bomb,''
\textit{Nature}, \textbf{238}, 211--212, 1972.

\bibitem{teukolsky1973}
S.~A. Teukolsky,
``Perturbations of a Rotating Black Hole. I.,''
\textit{Astrophys.\ J.}, \textbf{185}, 635--647, 1973.

\bibitem{detweiler1980}
S.~Detweiler,
``Klein-Gordon Equation and Rotating Black Holes,''
\textit{Phys.\ Rev.\ D}, \textbf{22}, 2323--2326, 1980.

\bibitem{dolan2007}
S.~R. Dolan,
``Instability of the Massive Klein-Gordon Field on the Kerr Spacetime,''
\textit{Phys.\ Rev.\ D}, \textbf{76}, 084001, 2007.

\bibitem{arvanitaki2015}
A.~Arvanitaki, M.~Baryakhtar, and X.~Huang,
``Discovering the QCD Axion with Black Holes and Gravitational Waves,''
\textit{Phys.\ Rev.\ D}, \textbf{91}, 084011, 2015.

\bibitem{ligo_sr2023}
R.~Abbott \textit{et al.}\ (LIGO Scientific, Virgo, and KAGRA
Collaborations),
``Searches for Gravitational Waves from Known Pulsars at Two Harmonics
in the Second and Third LIGO-Virgo Observing Runs,''
\textit{Astrophys.\ J.}, \textbf{935}, 1, 2022.

\bibitem{wald1984}
R.~M. Wald,
\textit{General Relativity},
University of Chicago Press, Chicago, 1984.

\end{thebibliography}

\end{document}
